\subsection{Doble BFS para Caminos Simples de Longitud "Más Corto + 1"}

\graybold{Preguntas guía}

\begin{itemize}
\item ¿Piden determinar si existe un camino SIMPLE (sin repetir vértices) de una longitud específica, distinta a la del camino más corto, entre dos nodos?
\item ¿La longitud pedida es exactamente uno más que la del camino más corto?
\item ¿El grafo no es ponderado, permitiendo usar BFS en vez de Dijkstra?
\end{itemize}

\graybold{¿Para qué sirve?} Determinar la existencia de un camino simple de longitud "camino más corto + 1" entre dos nodos de un grafo no ponderado, sin enumerar caminos explícitamente (exponencial), clasificando las aristas según los niveles BFS que conectan.

\graybold{Complejidad:} \bigO{n+m} (dos BFS y un recorrido de aristas).

\graybold{Fundamento teórico:} Sea $D$ la distancia más corta entre el origen (nodo 1) y el destino (nodo $n$), y $\text{dist1}(x)$ la distancia BFS desde el origen. En cualquier camino, cada arista cambia $\text{dist1}$ en $-1$, $0$, o $+1$ (propiedad de BFS en grafos no ponderados). Un camino de longitud $D+1$ necesita lograr un cambio neto de $D$ niveles en $D+1$ pasos; un conteo simple muestra que la ÚNICA forma posible es usar $D$ aristas "hacia adelante" ($+1$ nivel) y EXACTAMENTE UNA arista "horizontal" (mismo nivel) — cualquier arista "hacia atrás" haría inviable llegar a $D+1$ en solo $D+1$ pasos. Por lo tanto, basta buscar una arista horizontal $(u,v)$ con $\text{dist1}(u)=\text{dist1}(v)=k$ tal que alguno de los dos extremos tenga distancia más corta hacia el destino (calculada con una segunda BFS, desde $n$) igual a $D-k$: eso garantiza un camino "todo hacia adelante" desde ese extremo hasta $n$, mientras que el otro extremo se alcanza trivialmente desde el origen con el propio árbol BFS (también "todo hacia adelante"). Ambos segmentos son automáticamente simples y disjuntos entre sí (salvo en la arista horizontal), porque cada uno solo toca niveles estrictamente crecientes.

\graybold{Ejercicio aplicado}

En un grafo no ponderado de $n$ intersecciones y $m$ calles, con un camino más corto conocido de la intersección 1 a la $n$, determinar si existe una ruta SIMPLE (sin repetir intersecciones) de longitud exactamente una calle más larga que la ruta más corta.

\graybold{I/O:} n ≤ 10\textsuperscript{5}, m ≤ 2×10\textsuperscript{5} → lectura total con tokenización (patrón 2); dos BFS (desde 1 y desde $n$) sobre el mismo grafo. Verificado contra fuerza bruta (todos los caminos simples) en 385 grafos aleatorios pequeños, sin discrepancias, y con \textasciitilde{}0.34s en el peor caso de tamaño máximo.

\graybold{Código solución}

\begin{lstlisting}[language=Python]
import sys
from collections import deque

def bfs(start, adj, n):
    dist = [-1] * (n + 1)
    dist[start] = 0
    q = deque([start])
    while q:
        u = q.popleft()
        for v in adj[u]:
            if dist[v] == -1:
                dist[v] = dist[u] + 1
                q.append(v)
    return dist

def solve():
    data = sys.stdin.buffer.read().split()
    idx = 0
    n = int(data[idx]); m = int(data[idx + 1]); idx += 2
    adj = [[] for _ in range(n + 1)]
    edges = []
    for _ in range(m):
        a = int(data[idx]); b = int(data[idx + 1]); idx += 2
        adj[a].append(b)
        adj[b].append(a)
        edges.append((a, b))

    dist1 = bfs(1, adj, n)      # niveles BFS desde el origen
    distN = bfs(n, adj, n)      # distancias mas cortas hacia el destino
    D = dist1[n]

    for (u, v) in edges:
        if dist1[u] == dist1[v]:          # arista "horizontal": mismo nivel BFS desde 1
            k = dist1[u]
            # ¿alguno de los dos extremos tiene un camino "todo hacia adelante" hasta n?
            if distN[u] == D - k or distN[v] == D - k:
                print("possible")
                return
    print("impossible")

solve()
\end{lstlisting}
